\section{Related Work}
\label{sec:related}

\subsection{Solution space characterisation in redistricting}

The redistricting literature has characterised solution spaces primarily
through Markov chain Monte Carlo (MCMC) ensemble methods.
DeFord, Duchin, and Solomon~\citep{deford2019recombination} introduced
ReCom, a recombination Markov chain that generates large ensembles of
valid redistricting plans.  Their work demonstrated that the space of
compact, balanced plans is large and well-connected, a finding consistent
with our seed sweep results.

Fifield et al.~\citep{fifield2020} proved mixing bounds for redistricting
Markov chains, establishing that the solution space can be traversed
efficiently from any starting point.  Their empirical results on small
states ($k \leq 8$) showed that partisan outcomes vary substantially
across the solution space --- but this variation is driven by geographic
segregation, not by optimisation objective choices.

Autry et al.~\citep{autry2020} developed the multi-scale merger framework,
showing that ensemble diversity is highest for moderate $k$ values and
geographically complex states --- consistent with our finding that
GA ($k{=}14$) and NC ($k{=}14$) exhibit the highest seed variance.

\subsection{Sensitivity analysis for graph partitioning}

The graph partitioning literature has not systematically studied seed
sensitivity for METIS, though it is well known that METIS results vary
across seeds~\citep{karypis1998}.
Sanders and Schulz~\citep{sanders2013think} showed that running multiple
seeds and taking the best result improves quality by 2--5\% on benchmark
graphs.  Our sweep of 10,000 seeds per state characterises the full
distribution rather than just the minimum.

Pellegrini's SCOTCH partitioner~\citep{pellegrini2008scotch} uses a
different random number generator and shows similar seed sensitivity
profiles on planar graphs, suggesting that seed sensitivity is a property
of the problem class rather than the specific implementation.

\subsection{Partisan implications of redistricting algorithms}

Chen and Rodden~\citep{chen2013unintentional} showed that geographic
sorting of partisan preferences produces systematic partisan asymmetry
even in ``neutral'' algorithms.  This geographic effect dominates
seed-level variance in our results.

Stephanopoulos and McGhee~\citep{stephanopoulos2015partisan} introduced
the efficiency gap as a measure of partisan asymmetry.  We do not use
the efficiency gap directly, but our seat-share variance results are
related: states with high efficiency gap sensitivity tend to have high
seed variance.

\subsection{Hash-based seed generation}

The DIA's SHA-256 seed formula provides unpredictability (no party can
know the seed before the census release) and reproducibility (the seed
can be independently verified from public data).  Cryptographic hash
functions in randomised algorithm seeding have been studied by Bernstein
and Lange~\citep{bernstein2019}, who showed that domain separation
(the \texttt{DIA\_SEED\_V1} version string) is necessary to prevent
cross-application seed correlation.
